\chapter{Study Design and Ethics}
\label{ch:study}

\section{Classroom data only}
\label{sec:naturalistic}
The original plan staged the evidence: four weeks of baseline labs, then sessions in
which volunteers swapped seats at a set minute, copied, or pasted on cue, then mock
remote interviews. In practice a lab cannot be controlled that tightly. Any given
session could be a student's first, a re-attempt, on a lab PC or on their own laptop,
in a lab or a competition, and some students may well copy from a language model. The
students are a class being taught, not participants in an experiment, and asking them
to cheat on cue would change the very behaviour being measured.

The thesis therefore draws every conclusion from naturalistic recordings: whatever real
classes produce, as they happen, with no staged phase, no separate baseline weeks, no
interviews and no survey. This costs ground truth: no recorded session is known to
contain a switch, and none is known to be clean. The design below recovers what it can
without labels. Switches with a known time are made by splicing two real sessions
(\cref{sec:splicing}). Real sessions give a false-alarm rate that is an upper bound
(\cref{sec:protocol}). Device changes that students make on their own give a natural
experiment (\cref{sec:devices}). Signs the editor records give weak labels for copying
(\cref{sec:weak-labels}).

\section{Ethics and data protection}
\label{sec:ethics}
The pilot runs with oral institutional approval from the Dean, the Head of Department and
the course faculty; written confirmation is pending and will be attached as
Appendix~B. Every student sees a consent screen on first sign-in that says what is
recorded (key timing, edits, pastes, runs, focus; never the screen or camera), that a
flag only ever leads to a conversation, and that the teacher can watch live. The design
follows India's Digital Personal Data Protection Act 2023 and treats the EU AI Act's
high-risk requirements for education as a target: raw key events are deleted after 28
days and only window summaries are kept; a student can have all their data erased; the
analysis uses pseudonymous aggregates and never names a student. Splicing
(\cref{sec:splicing}) recombines summaries of two students' typing inside the analysis
only; a spliced session is never shown to a teacher or attached to either student.

\section{Setting and data}
\label{sec:pilot-data}
The editor has run in real IIIT Surat classes since 27 October 2026 (\cref{tab:pilot}).
Students use it for programming labs and class competitions in C and Python, on the
department's lab PCs or on their own laptops. Only pseudonymous aggregates were analysed
(\texttt{studentide.pilot}), and nothing is labelled.

\begin{table}[ht]
  \centering
  \caption{Pilot recordings, 27 Oct to 23 Nov 2026 (aggregates over all students, before the split was fixed).}
  \label{tab:pilot}
  \begin{tabular}{@{}lr@{}}
    \toprule
    Sessions (students) & 1,118 (281) \\
    Lab / competition / practice / learning lab & 443 / 590 / 50 / 35 \\
    C / Python & 608 / 456 \\
    Windows / Linux / Android / Mac / iOS & 470 / 444 / 111 / 84 / 9 \\
    Median session & 5 min, 108 key presses \\
    Events lost outside the 28-day key purge & 0 \\
    Sessions scored (baseline ready, some typing) & 157 \\
    Detector episodes per active hour, labs & 7.5 (process 4.2, pairs 2.7, identity 0.24) \\
    \quad on the student's usual / another device & 5.7 / 9.9 \\
    Study sessions with a paste / a big paste or paste after looking away & 248 / 117 of 760 \\
    \bottomrule
  \end{tabular}
\end{table}

Three lessons changed the plan. First, the recorder held up: a server restart during a
lab lost no events, because each browser kept and re-sent its batch. Second, phones send
no key events, so a student who codes on a phone has no rhythm at all, which is a fairness
limit for RQ4 and a reason for labs to require a laptop or PC. Third, the process channel
flagged far more often than on synthetic data, mostly through build order, which follows
the question more than the person, and flags rose on a device other than the student's
usual one. Since no takeover is known, these rates are an upper bound on false alarms;
they motivated re-tuning the detector on the development split (\cref{sec:split}) and a
device-aware baseline (\cref{sec:baseline}).

\paragraph{What counts as study data.} Learning-lab practice and teachers' ``View as
student'' previews are recorded but never analysed, edits that reset the editor to its
starter code are skipped, and sessions with no active window (30\,s with at least
20 code key presses) are left out so that idle tabs do not count as honest hours. Phone
sessions record no key events and so contribute only to the process channel; they are
reported separately.

\section{Baselines from the earliest minutes}
\label{sec:baseline}
Without baseline weeks, a student's baseline is simply their earliest $M$ active minutes
of typing, from whatever sessions those were in ($M = 20$ by default; the dev split
sweeps it). Sessions that contribute to the baseline are never scored. Since flags rose
on a student's less usual device, the baseline can be kept per device class (lab-PC-like,
personal-like, desktop app): a session is then scored against the student's baseline on
the same class, and the dev split decides whether this helps. Device classes come from a
browser fingerprint without its version: one shared by three or more students looks
like a lab PC. Raw keys of baseline sessions are purged after 28 days like any others;
the stored windows keep the baseline usable.

\section{Spliced takeovers}
\label{sec:splicing}
A spliced session takes a real, scored session of student $A$ and, from a time $\tau$,
continues it with the real typing of student $B$ from the same point of $B$'s session on
the same question, shifted to follow on. Every key, edit and run on either side was
really typed in class; only the join is made. The switch time is known by construction,
so detection, delay and the error of $\hat\tau$ can be measured on real typing. Rules
(\texttt{harness/splice.py}):
\begin{itemize}
  \item $B$ is another student on the same question and on the same side of the
    dev/test split, so the task, its stage and the split stay matched.
  \item A partner on the same device class gives the main condition; one on another
    device class gives a second, easier condition, since the device change is part of
    the signal.
  \item $\tau$ is drawn uniformly from multiples of 90\,s that leave at least
    three minutes, and two active windows, of each student.
  \item Splices are made from raw events when both sessions still have them (so the
    keystroke baselines can be scored), and from stored windows otherwise.
\end{itemize}

\begin{table}[ht]
  \centering
  \caption{Conditions in the classroom evaluation.}
  \label{tab:conditions}
  \begin{tabular}{@{}lll@{}}
    \toprule
    Condition & Built from & Answers \\
    \midrule
    Real session & a scored session as recorded & false alarms (upper bound) \\
    Splice, same device class & $A$ then $B$, same question & detection, delay (RQ1) \\
    Splice, other device class & $A$ then $B$, same question & detection with a device change \\
    Cross-device baseline & real sessions scored on the other class's baseline & transfer (RQ3) \\
    Copying signs & real sessions with weak signs & agreement (RQ2) \\
    \bottomrule
  \end{tabular}
\end{table}

What a splice cannot show is stated with every result. $B$ continues their own code,
not $A$'s, so line-level build order reflects $B$'s document, and a real helper would
first have to read $A$'s code, a pause that splices lack. Both make the process channel's
task easier than a real takeover, so splice detection is optimistic for that channel,
and the identity channels are also reported on their own. A splice also assumes that $A$
and $B$ are each a single author; if either copied or had help, the splice is a harder
case, not an easier one.

\section{Devices as a natural experiment}
\label{sec:devices}
Students were not assigned to devices; they used what they had. Many used both a lab PC
and their own laptop over the term, which allows a within-student comparison: the
distance between two sessions of the same student on the same device, on different
devices, and of two students on the same question. The same students are scored with a
baseline from one device class and from the other. Because device choice is not random
(students who bring laptops may differ), RQ3 compares within students only, and RQ4
reports device class as a group with that caveat.

\section{Weak labels for copying}
\label{sec:weak-labels}
The editor records signs that copying may have happened: a paste of 80 or more
characters, a paste within a minute of coming back to the window, text that appears
with no key press (an auto-typer), and, once teachers turn paste blocking on, each
attempt to paste text from outside the editor, which is blocked and logged. These are
signs, not labels: students also paste their own code. RQ2 asks whether transcription
measures agree with them; no accuracy is claimed.

\section{Held-out split: by student and by week}
\label{sec:split}
Two cuts make the test set. Thirty per cent of students form the test students, chosen
by a salted hash of their id, so the split does not depend on who else enrolled. In
addition, every session from 24 November 2026 onward is held out for all students, so
the test also covers later weeks the design never saw. Development (dev) data is
therefore the dev students' sessions before 24 November; detector re-tuning, feature
choices, the baseline length $M$, device-aware baselines and thresholds are decided on it
alone, and then frozen (\texttt{split.pilot\_split}, shared by the re-tuning work and
the harness). Test students' earliest minutes serve as their own baseline but are never
used for tuning.

\paragraph{Disclosure.} The first look at the pilot (\cref{tab:pilot}) reported
aggregates over all sessions to 23 November, test students included, before the split
was fixed. It informed the decision to re-tune and to make baselines device-aware, but
no parameter was fitted on it. Test-split sessions also feed class-level pieces of the
live review (feature weights, key-pair vocabulary); the frozen evaluation configuration
fits those on dev students only.

\section{Evaluation protocol}
\label{sec:protocol}
All research questions are answered by one evaluation harness
(\texttt{studentide.harness}, \cref{fig:harness}). It ran first on simulated typists,
whose results are kept as a reference; the classroom corpus
(\texttt{--db \dots{} --classroom}) replaces them with real sessions and splices.

\paragraph{Design on dev, report on test once.} All choices (the calibration source,
the process-channel window, the baseline length and whether it is device-aware, each
baseline method's alarm rule and CUSUM $k$) are made on dev students with 5-fold
cross-fitting, so that dev figures are also out-of-sample, and then frozen in a
configuration file. The test split is scored once per frozen configuration; every test
run is recorded in a ledger with its time, code version, configuration hash and corpus
fingerprint, and re-scoring the same corpus with the same configuration needs a written
reason, which is logged. Ablations are reported on dev, headline results on test.

\paragraph{Thresholds.} $h$ is set per channel for a target of 0.2 false alarms per
active hour of real sessions (split equally across channels). Calibration uses other
dev students' real sessions (\emph{holdout}), each student's own baseline sessions
scored leave-one-session-out (\emph{self}), or both (\emph{pooled}). Since some real
sessions may contain undetected cheating, calibrating on them sets $h$ slightly high,
which is the safe direction for a tool whose flags lead to conversations.

\section{Metrics}
For spliced sessions: the share with an alarm after $\tau$, the median and
90th-percentile delay, the median error of $\hat\tau$, and alarms before $\tau$. For real
sessions: false alarms per active hour and the share flagged at least once, both upper
bounds. Detection against false alarms is traced by sweeping the calibration target.
Precision at prevalence $\pi$ is $\mathrm{TPR}\,\pi/(\mathrm{TPR}\,\pi +
\mathrm{FPR}(1-\pi))$, with TPR from same-device splices and FPR from real sessions, at
$\pi = 1/100$ and $1/500$; using an upper-bound FPR makes this a lower bound on
precision. All intervals are 95\% percentile intervals from resampling students.

\begin{figure}[ht]
  \centering
  \resizebox{\linewidth}{!}{\tikzfig{evaluation-harness}}
  \caption{The evaluation harness: one protocol for simulated and classroom data.}
  \label{fig:harness}
\end{figure}

\section{Flags in teaching}
\label{sec:teacher-review}
No teacher is asked for extra time to review flags for the study. A flag is discussed when
it needs to be, as part of ordinary teaching: the teacher opens the evidence timeline and
talks with the student, and the flag leads to nothing else. The study does not stage or
script these conversations. After one, the teacher can note its outcome with one tap on
the moment (explained, own work; had help; copied; not discussed). These notes are
reported descriptively, by channel and context, as evidence of whether moments were worth
raising. They are not labels of cheating: a teacher's judgement after a short conversation
is itself uncertain, students who were never flagged are never asked, and the notes are
never used to tune the detector. Where no outcome is noted, the review interface is
assessed by its design (\cref{ch:system}).

\section{Threats to validity}
\begin{itemize}
  \item \emph{No ground truth.} Real sessions may hide real switches, so false alarms are
    overstated and precision understated; splices may differ from real takeovers
    (\cref{sec:splicing}).
  \item \emph{Self-selected devices.} Device effects are compared within students only.
  \item \emph{Little judgement of usefulness.} Only flags that teachers chose to discuss
    say anything about whether a flag was worth raising (\cref{sec:teacher-review}).
  \item \emph{One institution, one term, two languages.} Results may not transfer to
    other courses or to typing in other languages.
  \item \emph{Short baselines.} Median sessions are about five minutes, so 20 active
    minutes may span several sessions and days; the dev split tests shorter and longer.
  \item \emph{Raw-key purge.} After 28 days only stored windows remain, so keystroke
    baselines can be compared only on recent sessions; run the analysis within that
    window when it matters.
\end{itemize}
